%% kapa: exam=1 round=35 year=2024 period=2 form=A subject=accounting subject_ko=회계학 seq=10 number=50
%% source: Q-Net 감정평가사 기출문제 (cst00309) artlSeq=5213229 file=raw/1008/5213229/2교시 문제지 원본.pdf
%% conversion: kordoc 4.24.51 -> tools/gen1.py (auto); review=proofread-sonnet; flags=none
\begin{question}{50}
20×1년 초 설립된 (주)감평은 커피머신 1대를 이전(￦300)하면서 2년간 일정량의 원두를 공급(￦100)하기로 하는 계약을 체결하여 약속을 이행하고 현금 ￦400을 수령하였다. 이 계약이 고객과의 계약에서 생기는 수익의 기준을 모두 충족할 때 수익 인식 5단계 과정에 따라 순서대로 옳게 나열한 것은? (단, 거래가격의 변동 요소는 고려하지 않는다.)
\begin{framedtext}
ㄱ. 거래가격을 ￦400으로 산정 \par
ㄴ. 고객과의 계약에 해당하는지 식별 \par
ㄷ. 거래가격 ￦400을 커피머신 1대 이전에 대한 수행의무 1(￦300)과 2년간 원두공급에 대한 수행의무 2(￦100)에 배분 \par
ㄹ. 커피머신 1대 이전의 수행의무 1과 2년간 원두 공급의 수행의무 2로 수행의무 식별 \par
ㅁ. 수행의무 1(￦300)은 커피머신이 인도되는 시점에 수익을 인식하며, 수행의무 2(￦100)는 2년간 기간에 걸쳐 수익인식
\end{framedtext}
\begin{choices}
  \item ㄱ → ㄴ → ㄷ → ㄹ → ㅁ
  \item ㄴ → ㄱ → ㅁ → ㄷ → ㄹ
  \item ㄴ → ㄹ → ㄱ → ㄷ → ㅁ
  \item ㅁ → ㄷ → ㄱ → ㄴ → ㄹ
  \item ㅁ → ㄹ → ㄴ → ㄱ → ㄷ
\end{choices}
\end{question}
\answer{3}
